% ============================================================
%  EXAMEN TEMA 4: INTRODUCCIÓN AL ÁLGEBRA — POLINOMIOS
%  Curso de Introducción — 2do Año
%  Duración: 90 minutos
%  Temas: Expresiones algebraicas, términos semejantes,
%         ordenamiento, suma, resta, escalares y combinadas
%  Autor: Ing. Gabriel Astudillo
%  Nota: enunciado limpio (sin pistas ni desarrollo). Las
%  soluciones están en la "Clave de Respuestas" al final.
% ============================================================

\documentclass[12pt, letterpaper]{article}

% ── Paquetes ─────────────────────────────────────────────────

\usepackage{mathwizards-guia}
\mwguia{Examen — Tema 4: Suma y Resta de Polinomios}{Ing. Gabriel Astudillo $\cdot$ Curso 2do Año}

\begin{document}
% ============================================================

% ── Portada / Título ─────────────────────────────────────────
\mwportada{Examen del Tema 4}{Introducción al Álgebra}{Expresiones Algebraicas, Términos Semejantes y Polinomios}

\vspace{4pt}
\begin{cuadroinfo}
  \small
  \textbf{Nombre:} \rule{0.55\linewidth}{0.4pt} \quad \textbf{Fecha:} \rule{0.15\linewidth}{0.4pt} \\[6pt]
  \textbf{Tiempo disponible:} 90 minutos \quad$\bullet$\quad \textbf{Puntaje total:} 100 puntos \\[4pt]
  \textbf{Instrucciones:}
  \begin{enumerate}[leftmargin=*]
    \item Lee cada ejercicio con calma antes de resolverlo.
    \item Muestra \emph{todos} tus pasos y ordena los polinomios antes de operar.
    \item Recuerda: al \textbf{restar}, cambia el signo de \emph{todos} los términos del sustraendo.
    \item Solo se pueden sumar o restar \emph{términos semejantes}.
    \item No se permite el uso de calculadora.
    \item Escribe con letra clara y revisa tu examen antes de entregarlo.
  \end{enumerate}
\end{cuadroinfo}

\vspace{8pt}

% ============================================================
\section{Expresiones Algebraicas y sus Elementos (10 puntos)}
% ============================================================

\begin{enumerate}[leftmargin=*]
  \item \textbf{(5 pts)} Identifica el coeficiente, la variable (o variables) y el exponente (o exponentes) de cada término:
        \[
          a)\ -6x^5 \qquad b)\ 7a^3b^2 \qquad c)\ -\frac{4}{9}y \qquad d)\ m
        \]

  \item \textbf{(5 pts)} En la expresión $-4x^4 + 3x^2 - 8x + 6$, identifica:
        \begin{enumerate}[label=\alph*)]
          \item el coeficiente de cada término;
          \item el grado del polinomio;
          \item el término independiente.
        \end{enumerate}
\end{enumerate}

\newpage

% ============================================================
\section{Términos Algebraicos y Términos Semejantes (15 puntos)}
% ============================================================

\begin{enumerate}[leftmargin=*]
  \item \textbf{(5 pts)} Calcula el grado de cada término:
        \[
          a)\ 9x^4 \qquad b)\ -2a^3b^5 \qquad c)\ 15 \qquad d)\ 3xy^2z
        \]

  \item \textbf{(5 pts)} Reduce:
        \[
          a)\ 6x^2 - 2x^2 + 5x^2 - x^2 \qquad b)\ 4a^2b - 7a^2b + a^2b \qquad c)\ -3m + 5m + 2n - 8n + mn
        \]

  \item \textbf{(5 pts)} Reduce la expresión:
        \[
          7x^3 - 2x + 4x^3 + 5x - 9 + x
        \]
\end{enumerate}

\newpage

% ============================================================
\section{Ordenamiento de Polinomios (10 puntos)}
% ============================================================

\begin{enumerate}[leftmargin=*]
  \item \textbf{(5 pts)} Ordena el polinomio $6x - 3x^4 + 2x^2 + 1$ en forma descendente y en forma ascendente. Indica su grado.

  \item \textbf{(5 pts)} Ordena en forma descendente $4y^2 - 7 + 2y^5 + y^3 - y$ e indica su grado.
\end{enumerate}

\newpage

% ============================================================
\section{Suma de Polinomios (15 puntos)}
% ============================================================

\begin{enumerate}[leftmargin=*]
  \item \textbf{(5 pts)} Suma:
        \[
          (5x^2 - 3x + 4) + (-2x^2 + 7x - 1)
        \]

  \item \textbf{(5 pts)} Suma:
        \[
          (3a^2b - 2ab^2 + 5) + (a^2b + 4ab^2 - 3)
        \]

  \item \textbf{(5 pts)} Suma los tres polinomios:
        \[
          (2x^2 - x + 3) + (x^2 + 4x - 5) + (-3x^2 + 2x + 1)
        \]
\end{enumerate}

\newpage

% ============================================================
\section{Resta de Polinomios (15 puntos)}
% ============================================================

\begin{enumerate}[leftmargin=*]
  \item \textbf{(5 pts)} Resta:
        \[
          (7x^2 + 2x - 6) - (4x^2 - 5x + 3)
        \]

  \item \textbf{(5 pts)} Resta:
        \[
          (5a^3 - 2a^2 + 4a - 1) - (2a^3 + 3a^2 - a + 6)
        \]

  \item \textbf{(5 pts)} Resta:
        \[
          (4x^2 - 3xy + y^2) - (2x^2 + 5xy - 3y^2)
        \]
\end{enumerate}

\newpage

% ============================================================
\section{Multiplicación de un Polinomio por un Escalar (15 puntos)}
% ============================================================

\begin{enumerate}[leftmargin=*]
  \item \textbf{(5 pts)} Multiplica:
        \[
          a)\ 5(3x^2 - 2x + 4) \qquad b)\ -4(2x^3 + x - 7)
        \]

  \item \textbf{(5 pts)} Multiplica (escalares fraccionarios):
        \[
          a)\ \frac{1}{3}(9x^2 - 6x + 12) \qquad b)\ -\frac{1}{2}(8x - 10)
        \]

  \item \textbf{(5 pts)} Multiplica y reduce:
        \[
          -2(3a - 4b) + 5(a + 2b)
        \]
\end{enumerate}

\newpage

% ============================================================
\section{Suma y Resta de Polinomios con Escalares (20 puntos)}
% ============================================================

\begin{enumerate}[leftmargin=*]
  \item \textbf{(5 pts)} Calcula y reduce:
        \[
          3(2x^2 - x + 5) - 2(x^2 + 4x - 1)
        \]

  \item \textbf{(5 pts)} Calcula y reduce:
        \[
          -4(x^2 - 3x) + 2(3x^2 + x - 5)
        \]

  \item \textbf{(5 pts)} Calcula y reduce:
        \[
          \frac{2}{3}(6x^2 - 9x + 3) - \frac{1}{4}(8x^2 - 12x + 4)
        \]

  \item \textbf{(5 pts)} Verifica con $x = 2$ la igualdad:
        \[
          2(3x^2 - x + 4) - 3(x^2 + 2x - 1) = 3x^2 - 8x + 11
        \]
\end{enumerate}

\newpage

% ============================================================
\ifmwclaves
  \section{Clave de Respuestas}
  % ============================================================

  \subsection*{Sección 1: Expresiones Algebraicas y sus Elementos}

  \begin{enumerate}[leftmargin=*, label=\textbf{\arabic*.}]
    \item a) coeficiente $-6$, variable $x$, exponente $5$.
          \quad b) coeficiente $7$, variables $a$ y $b$, exponentes $3$ y $2$.
          \quad c) coeficiente $-\dfrac{4}{9}$, variable $y$, exponente $1$.
          \quad d) coeficiente $1$, variable $m$, exponente $1$.

    \item a) Coeficientes: $-4$ (en $x^4$), $3$ (en $x^2$), $-8$ (en $x$), $6$ (término independiente).
          \quad b) Grado $4$ (el mayor exponente). \quad c) Término independiente: $6$.
  \end{enumerate}

  \newpage

  \subsection*{Sección 2: Términos Algebraicos y Términos Semejantes}

  \begin{enumerate}[leftmargin=*, label=\textbf{\arabic*.}]
    \item a) grado $4$ \quad b) grado $3 + 5 = 8$ \quad c) grado $0$ \quad d) grado $1 + 2 + 1 = 4$

    \item a) $(6 - 2 + 5 - 1)x^2 = 8x^2$ \quad
          b) $(4 - 7 + 1)a^2b = -2a^2b$ \quad
          c) $(-3 + 5)m + (2 - 8)n + mn = 2m - 6n + mn$

    \item $(7 + 4)x^3 + (-2 + 5 + 1)x - 9 = 11x^3 + 4x - 9$
  \end{enumerate}

  \newpage

  \subsection*{Sección 3: Ordenamiento de Polinomios}

  \begin{enumerate}[leftmargin=*, label=\textbf{\arabic*.}]
    \item Descendente: $-3x^4 + 2x^2 + 6x + 1$. \quad Ascendente: $1 + 6x + 2x^2 - 3x^4$. \quad Grado: $4$.

    \item Descendente: $2y^5 + y^3 + 4y^2 - y - 7$. \quad Grado: $5$.
  \end{enumerate}

  \newpage

  \subsection*{Sección 4: Suma de Polinomios}

  \begin{enumerate}[leftmargin=*, label=\textbf{\arabic*.}]
    \item $(5 - 2)x^2 + (-3 + 7)x + (4 - 1) = 3x^2 + 4x + 3$

    \item $(3 + 1)a^2b + (-2 + 4)ab^2 + (5 - 3) = 4a^2b + 2ab^2 + 2$

    \item $(2 + 1 - 3)x^2 + (-1 + 4 + 2)x + (3 - 5 + 1) = 0x^2 + 5x - 1 = 5x - 1$
  \end{enumerate}

  \newpage

  \subsection*{Sección 5: Resta de Polinomios}

  \begin{enumerate}[leftmargin=*, label=\textbf{\arabic*.}]
    \item Cambiamos signos: $-4x^2 + 5x - 3$. \quad $(7 - 4)x^2 + (2 + 5)x + (-6 - 3) = 3x^2 + 7x - 9$

    \item Cambiamos signos: $-2a^3 - 3a^2 + a - 6$. \quad $(5-2)a^3 + (-2-3)a^2 + (4+1)a + (-1-6) = 3a^3 - 5a^2 + 5a - 7$

    \item Cambiamos signos: $-2x^2 - 5xy + 3y^2$. \quad $(4-2)x^2 + (-3-5)xy + (1+3)y^2 = 2x^2 - 8xy + 4y^2$
  \end{enumerate}

  \newpage

  \subsection*{Sección 6: Multiplicación por un Escalar}

  \begin{enumerate}[leftmargin=*, label=\textbf{\arabic*.}]
    \item a) $15x^2 - 10x + 20$ \quad b) $-8x^3 - 4x + 28$

    \item a) $\dfrac{1}{3}(9x^2) - \dfrac{1}{3}(6x) + \dfrac{1}{3}(12) = 3x^2 - 2x + 4$ \quad
          b) $-4x + 5$

    \item $-6a + 8b + 5a + 10b = (-6 + 5)a + (8 + 10)b = -a + 18b$
  \end{enumerate}

  \newpage

  \subsection*{Sección 7: Suma y Resta de Polinomios con Escalares}

  \begin{enumerate}[leftmargin=*, label=\textbf{\arabic*.}]
    \item $6x^2 - 3x + 15 - 2x^2 - 8x + 2 = 4x^2 - 11x + 17$

    \item $-4x^2 + 12x + 6x^2 + 2x - 10 = 2x^2 + 14x - 10$

    \item $\dfrac{2}{3}(6x^2 - 9x + 3) = 4x^2 - 6x + 2$ \quad y \quad $\dfrac{1}{4}(8x^2 - 12x + 4) = 2x^2 - 3x + 1$.
          Resultado: $(4x^2 - 6x + 2) - (2x^2 - 3x + 1) = 2x^2 - 3x + 1$.

    \item Lado izquierdo: $2(3x^2 - x + 4) - 3(x^2 + 2x - 1) = 6x^2 - 2x + 8 - 3x^2 - 6x + 3 = 3x^2 - 8x + 11$.
          Ambas expresiones coinciden para todo valor de $x$ (es una identidad).
          Con $x = 2$: izquierda $= 2(12 - 2 + 4) - 3(4 + 4 - 1) = 2(14) - 3(7) = 28 - 21 = 7$; derecha $= 3(4) - 8(2) + 11 = 12 - 16 + 11 = 7$. $\checkmark$
  \end{enumerate}

\fi
\end{document}
